% ID: SEL-D0078-OA-E03
% ESTADO: cotejada. Segunda pasada visual y borrador compilado revisados.
% Fuente: Selectividad/sel_2008_matematicas_aplicadas/Mat CCSS 6 - Examen y criterios de correción.pdf
% SHA-256: 256ad5f6dc324a5bd081279a0455d3ea64e9feb86c8be52341faa83024bc3b5b
% Página: 1; opción A; ejercicio original 3.

\noindent\underline{Parte I}

\textbf{(2 puntos)} Ana y Blas deciden jugar con un dado de la siguiente forma:

``Ana lanza el dado y, si saca un 6, gana y se acaba el juego. En caso contrario lanza Blas, que gana si saca un 2 o un 3, y también se acaba el juego. De no ocurrir esto, la partida se acaba sin ganador.

Halle la probabilidad de los siguientes sucesos: ``gana Ana'', ``gana Blas'', ``ninguno gana''.
\par\noindent\begin{minipage}{\linewidth}
\noindent\underline{Parte II}

\textbf{(2 puntos)} En una muestra representativa de 1200 residentes de una ciudad, 450 utilizan habitualmente el transporte público. Obtenga el intervalo de confianza, al $90\%$, de la proporción de residentes en la ciudad que utilizan habitualmente el transporte público.
\end{minipage}\par
